\documentclass{article}
\usepackage{amsmath,amssymb}
\pagestyle{plain}
\begin{document}
\section{Limits}
For every $\varepsilon>0$ there is $N\in\mathbb{N}$ with $|a_n-L|<\varepsilon$ for all $n\ge N$, so
\begin{equation}
\lim_{n\to\infty}\left(1+\frac{1}{n}\right)^{n}=e .
\end{equation}
A nested fraction: $\frac{1}{1+\frac{1}{x}}=\frac{x}{x+1}$ and a root $\sqrt{x^2+y^2}$.
\[
\frac{d}{dx}\sin^2 x = 2\sin x\cos x, \qquad \int_{-\infty}^{\infty} e^{-x^2}\,dx=\sqrt{\pi}
\]
The binomial theorem $(a+b)^n=\sum_{k=0}^{n}\binom{n}{k}a^k b^{n-k}$ and
\[
\prod_{i=1}^{n} x_i \le \left(\frac{x_1+\cdots+x_n}{n}\right)^{n}
\]
Vectors: $\vec{v}\cdot\vec{w}=\|v\|\,\|w\|\cos\theta$ and $f'(x)=\tilde{g}(x)$.
\end{document}
